% ============================================================================
%  CHƯƠNG 3. CÀI ĐẶT, THỰC NGHIỆM VÀ ỨNG DỤNG
% ============================================================================
\chapter{Cài đặt, thực nghiệm và ứng dụng}\label{chap:3}

Chương này kiểm nghiệm các kết quả lý thuyết của Chương~\ref{chap:2} bằng thực nghiệm và trả lời bốn câu hỏi nghiên cứu nêu ở Mở đầu. \Cref{sec:impl} mô tả phần cài đặt và các phép kiểm chứng số. \Cref{sec:exp2d} dùng dữ liệu hai chiều để quan sát trực tiếp biên quyết định. \Cref{sec:benchmark} tái hiện một thí nghiệm của bài báo gốc, so sánh ba biến thể PWL-SVM với sáu phương pháp khác trên mười bộ dữ liệu chuẩn và khảo sát ảnh hưởng của số đặc trưng cùng quy tắc chọn tham số. \Cref{sec:cost} đo chi phí huấn luyện, dự đoán và bộ nhớ. \Cref{sec:application} áp dụng PWL-SVM vào bài toán dự báo vỡ nợ thẻ tín dụng: so sánh với thẻ điểm truyền thống và các mô hình phi tuyến, phân tích khả năng diễn giải, và bổ sung ràng buộc đơn điệu. \Cref{sec:discussion} tổng hợp và thảo luận các kết quả.

% ----------------------------------------------------------------------------
\section{Cài đặt}\label{sec:impl}

Toàn bộ thực nghiệm được cài đặt bằng Python~3 với các thư viện NumPy~\cite{harris2020}, SciPy~\cite{virtanen2020} và scikit-learn~\cite{pedregosa2011}. Mã nguồn gồm một mô-đun chính, \texttt{pwlsvm.py}, với hai lớp tuân theo giao diện \texttt{fit}/\texttt{predict} của scikit-learn; nhờ vậy có thể dùng trực tiếp các công cụ kiểm định chéo và tìm kiếm tham số của thư viện này.
\begin{itemize}
\item Lớp \texttt{PWLFeatures} sinh tham số và tính ánh xạ đặc trưng~\eqref{eq:feature-map} cho ba dạng cộng tính, HH và GHH, theo đúng quy trình ở Mục~\ref{sec:maps}. Ngoài quy tắc của bài báo gốc, lớp này có hai tùy chọn: đặt nút của dạng cộng tính theo phân vị (\cref{rem:quantile}) và chuẩn hóa các siêu phẳng ngẫu nhiên về $\|p\|_2=1$ (\cref{rem:scale}).
\item Lớp \texttt{PWLSVM} ghép ba bước: chuẩn hóa từng thành phần về $[0,1]$ (giá trị nhỏ nhất, lớn nhất được học trên tập huấn luyện), tính $\varphi(x)$, và giải bài toán tuyến tính trên $\varphi(x)$. Với PWL-C-SVM có hai bộ giải. LIBSVM~\cite{chang2011} giải đúng bài toán đối ngẫu~\eqref{eq:pwl-c-dual} với hệ số chặn không bị phạt; nó được dùng cho các thí nghiệm nhỏ ở \cref{sec:exp2d} và \cref{sec:cosexp}. LIBLINEAR~\cite{fan2008} giải bài toán trên đặc trưng tường minh bằng phương pháp hạ theo tọa độ trên bài toán đối ngẫu, nhanh hơn nhiều khi $N$ hoặc $\gamma$ lớn; hệ số chặn được đưa vào như trọng số của một đặc trưng hằng bằng $10$ nên bị chính quy hóa nhẹ, với phạt tương đương $w_0^2/200$. LIBLINEAR được dùng cho mọi thí nghiệm còn lại. Với PWL-LS-SVM, dạng gốc được giải trực tiếp theo \cref{prop:ridge}.
\end{itemize}
Mô-đun thứ hai, \texttt{pwlsvm\_monotone.py}, bổ sung ràng buộc đơn điệu cho PWL-LS-SVM cộng tính (Mục~\ref{sec:monotone}). Các phương pháp đối chứng --- SVM tuyến tính, SVM và LS-SVM hạt nhân Gauss, SVM hạt nhân giao, kNN, mạng nơ-ron một lớp ẩn ReLU, và trong phần ứng dụng thêm hồi quy logistic, thẻ điểm logistic trên biến rời rạc hóa và tổ hợp cây tăng cường --- được đặt trong cùng một khung, để mọi phương pháp dùng chung cách chuẩn hóa, cách chia dữ liệu và cách chọn tham số. Quy trình huấn luyện và dự đoán của PWL-SVM được tóm tắt trong \cref{alg:pwlsvm}.

\begin{algorithm}[t]
\caption{Huấn luyện và dự đoán với PWL-SVM}\label{alg:pwlsvm}
\begin{algorithmic}[1]
\Require dữ liệu $\{(x_k,y_k)\}_{k=1}^{N}$; dạng ánh xạ (cộng tính, HH hoặc GHH); số đặc trưng trên mỗi chiều $M/n$; hằng số $\gamma$
\Ensure trọng số $w\in\R^M$, hệ số chặn $w_0$, tham số của ánh xạ và của phép chuẩn hóa
\State Chuẩn hóa từng thành phần của $x_k$ về $[0,1]$; lưu giá trị nhỏ nhất và lớn nhất
\If{dạng cộng tính}
  \State đặt $S=M/n$ và các nút $q=s/S$ ($s=1,\dots,S-1$), hoặc các phân vị mức $s/S$, trên mỗi trục
\Else
  \For{$m=n+1,\dots,M$}
    \State sinh $1$ siêu phẳng (HH) hoặc $n$ siêu phẳng (GHH) theo~\eqref{eq:random-plane}
  \EndFor
\EndIf
\State Lập ma trận $\Phi\in\R^{N\times M}$ có hàng thứ $k$ là $\varphi(x_k)^\T$, theo~\eqref{eq:feature-map}
\State Giải~\eqref{eq:pwl-c-dual} rồi tính $w=\sum_k\alpha_ky_k\varphi(x_k)$ (PWL-C-SVM), hoặc giải~\eqref{eq:ridge} (PWL-LS-SVM)
\State \textbf{Dự đoán} điểm mới $x$: chuẩn hóa $x$ bằng các giá trị đã lưu, trả về $\sgn\bigl(w^\T\varphi(x)+w_0\bigr)$
\end{algorithmic}
\end{algorithm}

\paragraph{Kiểm chứng tính đúng đắn.} Trước khi dùng mã cho thực nghiệm, đề án kiểm tra bằng số hai hệ thức của Chương~\ref{chap:2}, trên dữ liệu hai mặt trăng ($N=400$, $M=20$, $\gamma=10$) với cả ba dạng ánh xạ.
\begin{enumerate}[label=(\roman*)]
\item Với PWL-LS-SVM, nghiệm của dạng gốc (hồi quy ridge, \cref{prop:ridge}) và nghiệm của hệ đối ngẫu~\eqref{eq:pwl-ls-dual} phải trùng nhau. Sai khác lớn nhất giữa hai bộ tham số $(w,w_0)$ là $7\cdot10^{-12}$; trên $400$ điểm kiểm tra, giá trị hàm quyết định của hai lời giải sai khác không quá $2\cdot10^{-12}$. Hai cách giải cho cùng một nghiệm, chỉ khác nhau ở sai số làm tròn.
\item Với PWL-C-SVM, khôi phục $w=\sum_k\alpha_ky_k\varphi(x_k)$ từ nghiệm đối ngẫu với hạt nhân hữu hạn hạng~\eqref{eq:kernel} cho vectơ trọng số trùng với vectơ thu được khi giải trực tiếp trên $\varphi(x)$ (sai khác dưới $2\cdot10^{-12}$); hai bộ phân loại dự đoán giống hệt nhau trên toàn bộ tập kiểm tra.
\end{enumerate}
Ngoài ra, ví dụ hai mặt trăng ở Mục~\ref{sec:maps} tái hiện đúng hiện tượng mà bài báo gốc mô tả: bốn đặc trưng cộng tính cho một biên ba đoạn với các điểm gãy nằm trên hai đường $x(2)=\tfrac13$ và $x(2)=\tfrac23$.

% ----------------------------------------------------------------------------
\section{Thực nghiệm minh họa trên dữ liệu hai chiều}\label{sec:exp2d}

Với dữ liệu hai chiều, ta nhìn được trực tiếp biên quyết định, nên đây là nơi thích hợp để kiểm tra các nhận định định tính ở Chương~\ref{chap:2}. Đề án dùng bốn bộ dữ liệu tổng hợp có độ khó tăng dần (\cref{fig:demo2d}): hai mặt trăng (nhiễu Gauss với độ lệch chuẩn $0{,}15$), hai vòng tròn đồng tâm (tỉ số bán kính $0{,}5$, nhiễu $0{,}08$), bàn cờ $4\times4$ với các điểm phân bố đều, và hai xoắn ốc lồng nhau ($1{,}6$ vòng, nhiễu $0{,}06$). Mỗi bộ gồm $1000$ điểm, chia đôi thành $500$ điểm huấn luyện và $500$ điểm kiểm tra như trong bài báo gốc.

Năm phương pháp được so sánh: SVM tuyến tính; PWL-C-SVM với ánh xạ cộng tính ($M/n=10$), HH ($M/n=50$, tức $98$ siêu phẳng) và GHH ($M/n=25$, tức $48$ đặc trưng với tổng cộng $96$ siêu phẳng); và SVM hạt nhân Gauss. Hai cấu hình HH và GHH được chọn để có số siêu phẳng xấp xỉ nhau. Hằng số $\gamma$ được chọn trong tập $\{2^{-3},2^{-1},\dots,2^{11}\}$ bằng kiểm định chéo $5$ phần trên tập huấn luyện; với SVM hạt nhân Gauss, tham số $1/\sigma^2$ được chọn đồng thời trong $\{2^{-3},2^{-1},\dots,2^{9}\}$. Toàn bộ quy trình được lặp lại $10$ lần với dữ liệu và tham số ngẫu nhiên khác nhau; \cref{tab:exp2d} báo cáo trung bình và độ lệch chuẩn.

\begin{figure}[p]
\centering
\includegraphics[width=\linewidth]{fig_demo_2d.pdf}
\caption[Biên quyết định của năm phương pháp trên dữ liệu hai chiều]{Biên quyết định của năm phương pháp trên bốn bộ dữ liệu hai chiều (lần lặp đầu tiên). Chấm xanh và dấu nhân cam là các điểm kiểm tra của hai lớp; con số ở góc là độ chính xác trên $500$ điểm kiểm tra. Biên của PWL-SVM là các đường gấp khúc, biên của SVM hạt nhân Gauss là đường trơn.}
\label{fig:demo2d}
\end{figure}

\begin{table}[t]
\centering
\caption[Độ chính xác trên dữ liệu hai chiều]{Độ chính xác trên tập kiểm tra (\%, trung bình $\pm$ độ lệch chuẩn qua $10$ lần lặp) trên dữ liệu hai chiều; in đậm là giá trị trung bình cao nhất mỗi hàng.}
\label{tab:exp2d}
\small
\renewcommand{\arraystretch}{1.2}
\setlength{\tabcolsep}{4pt}
\begin{tabular}{@{}l c c c c c@{}}
\toprule
\textbf{Dữ liệu} & \textbf{SVM tuyến tính} & \textbf{PWL cộng tính} & \textbf{PWL-HH} & \textbf{PWL-GHH} & \textbf{SVM-RBF}\\
\midrule
Hai mặt trăng & $87{,}3\pm2{,}0$ & $98{,}8\pm0{,}7$ & $98{,}7\pm0{,}7$ & $98{,}6\pm0{,}8$ & $\mathbf{98{,}9}\pm0{,}6$\\
Hai vòng tròn & $52{,}7\pm3{,}5$ & $\mathbf{99{,}9}\pm0{,}1$ & $99{,}8\pm0{,}2$ & $99{,}8\pm0{,}2$ & $99{,}8\pm0{,}2$\\
Bàn cờ $4\times4$ & $50{,}0\pm2{,}4$ & $55{,}5\pm3{,}4$ & $90{,}6\pm2{,}2$ & $85{,}7\pm2{,}5$ & $\mathbf{94{,}8}\pm1{,}2$\\
Hai xoắn ốc & $64{,}8\pm2{,}7$ & $89{,}2\pm1{,}5$ & $98{,}2\pm0{,}6$ & $96{,}7\pm1{,}4$ & $\mathbf{99{,}2}\pm0{,}4$\\
\bottomrule
\end{tabular}
\end{table}

Kết quả ở \cref{tab:exp2d} và \cref{fig:demo2d} phù hợp với phân tích ở Mục~\ref{sec:maps}. SVM tuyến tính thất bại trên cả bốn bộ; trên hai vòng tròn và bàn cờ, độ chính xác của nó xấp xỉ $50\%$, tức không hơn đoán ngẫu nhiên. Cả ba dạng PWL-SVM đều cải thiện rõ rệt, và trên hai mặt trăng, hai vòng tròn, chúng đạt độ chính xác ngang SVM hạt nhân Gauss, với sai khác nhỏ hơn độ lệch chuẩn.

Ánh xạ cộng tính cho kết quả gần như tuyệt đối trên hai vòng tròn ($99{,}9\%$) nhưng gần như thất bại trên bàn cờ ($55{,}5\%$). Nguyên nhân không nằm ở số đặc trưng mà ở cấu trúc của mô hình. Biên của hai vòng tròn là đường $x(1)^2+x(2)^2=r^2$, với vế trái là tổng của hai hàm một biến, nên một mô hình cộng tính xấp xỉ được dễ dàng. Ngược lại, mệnh đề sau cho thấy không mô hình cộng tính nào xử lý được cấu trúc kiểu bàn cờ.

\begin{proposition}\label{prop:xor}
Cho $a\neq b$ và $c\neq d$. Giả sử hai điểm $(a,c)$ và $(b,d)$ thuộc lớp $+1$, còn hai điểm $(a,d)$ và $(b,c)$ thuộc lớp $-1$. Khi đó không tồn tại hai hàm $g_1,g_2:\R\to\R$ nào sao cho hàm cộng tính $f(x)=g_1\bigl(x(1)\bigr)+g_2\bigl(x(2)\bigr)$ phân loại đúng cả bốn điểm.
\end{proposition}

\begin{proof}
Nếu có, thì $f(a,c)+f(b,d)>0$ và $f(a,d)+f(b,c)<0$. Nhưng cả hai tổng đều bằng $g_1(a)+g_1(b)+g_2(c)+g_2(d)$, mâu thuẫn.
\end{proof}

Bàn cờ $4\times4$ chứa rất nhiều bộ bốn điểm như vậy, nên mọi bộ phân loại cộng tính, dù dùng bao nhiêu đặc trưng, đều mắc nhiều lỗi. Kết luận này áp dụng cả cho SVM hạt nhân giao, vì theo Mục~\ref{sec:relations} nó cũng là một mô hình cộng tính; điều đó giải thích vì sao trong bài báo gốc, SVM hạt nhân giao chỉ đạt $48{,}8\%$ trên dữ liệu bàn cờ~\cite[Bảng~1]{huang2013}.

Hai dạng HH và GHH, vốn dùng các siêu phẳng có hướng tùy ý, xử lý được cả bàn cờ lẫn xoắn ốc. Trên hai bộ khó này, chúng vẫn kém SVM hạt nhân Gauss vài điểm phần trăm ($90{,}6\%$ và $85{,}7\%$ so với $94{,}8\%$ trên bàn cờ), và biên của chúng gấp khúc thay vì trơn. Với cùng số siêu phẳng, GHH kém HH một chút: mỗi đặc trưng GHH linh hoạt hơn, nhưng số trọng số được học chỉ bằng một nửa ($50$ so với $100$), trong khi các siêu phẳng ngẫu nhiên không được điều chỉnh theo dữ liệu. Nói cách khác, khả năng biểu diễn tổng quát của GHH (\cref{cor:boundary}) không tự động chuyển thành độ chính xác khi tham số được sinh ngẫu nhiên.

Một nhận xét về quy trình thực nghiệm: trên bàn cờ, hằng số $\gamma$ được chọn thường nằm ở đầu trên của lưới ($2^{9}$ hoặc $2^{11}$), và thời gian chọn tham số của PWL-HH và PWL-GHH (khoảng $9$ và $18$ giây mỗi lần lặp) lớn hơn nhiều so với SVM hạt nhân Gauss (dưới $0{,}4$ giây), vì LIBSVM hội tụ chậm khi $\gamma$ lớn. Đây là lý do các thí nghiệm trên dữ liệu chuẩn ở \cref{sec:benchmark} dùng bộ giải LIBLINEAR.

% ----------------------------------------------------------------------------

% ----------------------------------------------------------------------------
\section{Thực nghiệm trên dữ liệu chuẩn}\label{sec:benchmark}
% ----------------------------------------------------------------------------
%  3.3  Thực nghiệm trên dữ liệu chuẩn
% ----------------------------------------------------------------------------
\subsection{Tái hiện thí nghiệm trên dữ liệu Cosexp}\label{sec:cosexp}

Bộ dữ liệu Cosexp của hộp công cụ SVM-KM~\cite{canu2005} được bài báo gốc dùng để minh họa ảnh hưởng của số đặc trưng $M$. Các điểm được lấy đều trên miền $[0,5]\times[-1,1]$ và được gán nhãn theo dấu của $\cos\bigl(0{,}5\,e^{0{,}7x(1)}\bigr)-x(2)$; biên thật là đường cong $x(2)=\cos\bigl(0{,}5\,e^{0{,}7x(1)}\bigr)$, dao động ngày càng nhanh khi $x(1)$ tăng (nét đứt ở \cref{fig:cosexp}a,b). Theo đúng thiết lập của~\cite{huang2013}, ta dùng $500$ điểm huấn luyện, $500$ điểm kiểm tra và PWL-C-SVM với ánh xạ cộng tính, cho $M$ chạy từ $2$ đến $40$. Khác với bài báo, mỗi cấu hình được lặp lại trên $10$ bộ dữ liệu sinh độc lập, và ánh xạ HH được chạy song song để so sánh.

\begin{figure}[t]
\centering
\includegraphics[width=\linewidth]{fig_cosexp_M.pdf}
\caption[Tái hiện thí nghiệm Cosexp: ảnh hưởng của số đặc trưng]{Tái hiện thí nghiệm Cosexp (Hình~3 của~\cite{huang2013}), $10$ lần lặp. (a), (b) Biên của PWL-C-SVM cộng tính với $M=12$ và $M=40$ (nét liền) so với biên thật (nét đứt). (c) Độ chính xác trên tập kiểm tra theo $M$: đường là trung bình, dải màu là $\pm$ một độ lệch chuẩn; đường chấm là SVM hạt nhân Gauss. (d) Thời gian huấn luyện (trung vị) với $\gamma$ đã chọn.}
\label{fig:cosexp}
\end{figure}

\Cref{fig:cosexp} tái hiện đúng xu hướng mà bài báo gốc báo cáo. Với $M=2$, bộ phân loại chỉ là tuyến tính và đạt $72{,}7\pm2{,}0\%$ (bài báo: $76{,}83\%$ với một lần chạy). Khi $M$ tăng, biên uốn theo đường cong thật ngày càng sát (\cref{fig:cosexp}a,b); độ chính xác đạt $90{,}0\%$ tại $M=20=10n$ và ổn định quanh $94\%$ từ $M=24$, ngang với SVM hạt nhân Gauss ($94{,}2\pm0{,}9\%$). Ở $M$ nhỏ, đường cong không đơn điệu: mỗi giá trị của $M$ ứng với một lưới nút hoàn toàn khác, và độ chính xác phụ thuộc vào việc các nút có rơi gần những chỗ biên đổi hướng hay không. Thời gian huấn luyện tăng theo $M$ nhưng vẫn dưới $0{,}1$ giây (\cref{fig:cosexp}d).

Đáng chú ý, trên dữ liệu này ánh xạ HH kém hẳn ánh xạ cộng tính ($89{,}1\pm1{,}4\%$ tại $M=40$), dù về lý thuyết nó linh hoạt hơn. Lý do là nhãn của Cosexp được xác định bởi hàm $\cos\bigl(0{,}5\,e^{0{,}7x(1)}\bigr)-x(2)$, một hàm cộng tính: ánh xạ cộng tính khớp đúng cấu trúc đó, còn các siêu phẳng ngẫu nhiên của HH phân tán khả năng biểu diễn vào những hướng không cần thiết. Kết hợp với Mục~\ref{sec:exp2d}, ta rút ra một quy tắc thực hành: nên thử ánh xạ cộng tính trước, vì nó rẻ, dễ diễn giải và đủ tốt khi tương tác giữa các biến yếu; chỉ chuyển sang HH khi có dấu hiệu tương tác mạnh, như trên bàn cờ. Đây cũng là thứ tự mà bài báo gốc khuyến nghị~\cite[Mục~4]{huang2013}.

% Generated by code/fill_numbers.py from chuong3_bench_b.tpl -- edit the template, not this file.
% ----------------------------------------------------------------------------
\subsection{Mười bộ dữ liệu chuẩn}\label{sec:bench10}

\paragraph{Dữ liệu và giao thức.} Đề án dùng mười bộ dữ liệu phân loại nhị phân (\cref{tab:datasets}). Bảy bộ trong số đó có mặt trong Bảng~1 của bài báo gốc (Haberman, Transfusion, Pima, Breast, Magic, Parkinsons, Ionosphere); ba bộ còn lại (Banknote, Spambase, Sonar) được thêm vào để bao quát thêm các kích thước dữ liệu và số chiều khác nhau. Mỗi bộ được chia ngẫu nhiên có phân tầng thành hai nửa huấn luyện và kiểm tra, lặp lại $10$ lần; với Magic, tập huấn luyện gồm $2.000$ điểm như trong bài báo gốc, phần còn lại dùng để kiểm tra; với Spambase, do chi phí tinh chỉnh các mô hình hạt nhân lớn, chỉ lặp lại $5$ lần. Trên mỗi tập huấn luyện, tham số được chọn bằng kiểm định chéo $10$ phần: $\gamma\in\{2^{-5},2^{-3},\dots,2^{11}\}$ cho mọi mô hình SVM, tham số $1/\sigma^2$ của hạt nhân Gauss trên lưới $\{2^{-15},2^{-13},\dots,2^{3}\}$ theo khuyến nghị của~\cite{hsu2003}, và hệ số phạt $L_2$ của mạng nơ-ron trong $\{10^{-4},10^{-3},\dots,1\}$. Chín phương pháp được so sánh: SVM tuyến tính; SVM và LS-SVM hạt nhân Gauss; kNN với $k=1$ như trong bài báo; Ik-SVM; mạng nơ-ron một lớp ẩn gồm $D=9n$ nơ-ron ReLU, bằng đúng số đặc trưng bản lề của PWL-SVM dạng HH, được huấn luyện bằng L-BFGS~\cite{liu1989}; và ba biến thể PWL-SVM với $M/n=10$ theo đúng quy tắc của bài báo: PWL-C-SVM và PWL-LS-SVM với ánh xạ cộng tính lưới đều, và PWL-C-SVM với ánh xạ HH chuẩn hóa $|p(1)|=1$.

\begin{table}[t]
\centering
\caption[Mười bộ dữ liệu chuẩn]{Mười bộ dữ liệu chuẩn, sắp theo số chiều $n$ (sau khi bỏ các cột hằng). $N$ là số mẫu, cột cuối là số mẫu huấn luyện/kiểm tra trong mỗi lần chia.}
\label{tab:datasets}
\small
\renewcommand{\arraystretch}{1.15}
\setlength{\tabcolsep}{4.5pt}
\begin{tabular}{@{}l l r r r r@{}}
\toprule
\textbf{Tên} & \textbf{Nội dung} & $N$ & $n$ & \begin{tabular}[b]{@{}c@{}}\textbf{\% lớp}\\ \textbf{$+1$}\end{tabular} & \begin{tabular}[b]{@{}c@{}}\textbf{Huấn luyện/}\\ \textbf{kiểm tra}\end{tabular}\\
\midrule
Haberman & sống sót sau phẫu thuật ung thư vú & 306 & 3 & 26{,}5 & 153/153\\
Transfusion & hiến máu tình nguyện & 748 & 4 & 23{,}8 & 374/374\\
Banknote & xác thực tiền giấy & 1.372 & 4 & 44{,}5 & 686/686\\
Pima & bệnh tiểu đường & 768 & 8 & 34{,}9 & 384/384\\
Breast & ung thư vú Wisconsin (bản gốc) & 699 & 9 & 34{,}5 & 349/350\\
Magic & kính thiên văn tia gamma & 19.020 & 10 & 35{,}2 & 2.000/17.020\\
Parkinsons & bệnh Parkinson (giọng nói) & 195 & 22 & 75{,}4 & 97/98\\
Ionosphere & tín hiệu radar tầng điện li & 351 & 33 & 64{,}1 & 175/176\\
Spambase & thư rác & 4.601 & 57 & 39{,}4 & 2.300/2.301\\
Sonar & tín hiệu sonar (đá hay mìn) & 208 & 60 & 53{,}4 & 104/104\\
\bottomrule
\end{tabular}
\end{table}


\begin{table}[t]
\centering
\caption[Độ chính xác trên mười bộ dữ liệu chuẩn]{Độ chính xác trung bình trên tập kiểm tra (\%) qua $10$ lần chia ngẫu nhiên ($5$ lần với Spambase). In đậm: giá trị cao nhất mỗi hàng. Dòng cuối là hạng trung bình trên mười bộ dữ liệu ($1$ là tốt nhất). Độ lệch chuẩn được cho ở \cref{tab:appendix-bench} (Phụ lục~B).}
\label{tab:bench}
\footnotesize
\setlength{\tabcolsep}{3pt}
\renewcommand{\arraystretch}{1.18}
\resizebox{\linewidth}{!}{%
\begin{tabular}{@{}l r r r r r r r r r@{}}
\toprule
\textbf{Dữ liệu} & \begin{tabular}[b]{@{}c@{}}SVM\\tuyến tính\end{tabular} & \begin{tabular}[b]{@{}c@{}}SVM-\\RBF\end{tabular} & \begin{tabular}[b]{@{}c@{}}LS-SVM-\\RBF\end{tabular} & \begin{tabular}[b]{@{}c@{}}kNN\\$(k=1)$\end{tabular} & \begin{tabular}[b]{@{}c@{}}Ik-\\SVM\end{tabular} & \begin{tabular}[b]{@{}c@{}}MLP-\\ReLU\end{tabular} & \begin{tabular}[b]{@{}c@{}}PWL-C\\cộng tính\end{tabular} & \begin{tabular}[b]{@{}c@{}}PWL-LS\\cộng tính\end{tabular} & \begin{tabular}[b]{@{}c@{}}PWL-C\\HH\end{tabular}\\
\midrule
Haberman & 72{,}5 & 73{,}3 & 73{,}4 & 65{,}7 & 72{,}0 & 72{,}4 & 73{,}2 & 72{,}6 & \textbf{73{,}5}\\
Transfusion & 76{,}5 & 77{,}6 & 77{,}7 & 68{,}4 & 76{,}7 & \textbf{78{,}6} & 76{,}3 & 77{,}4 & 76{,}3\\
Banknote & 98{,}6 & \textbf{100{,}0} & \textbf{100{,}0} & 99{,}8 & 99{,}9 & 99{,}9 & 99{,}9 & 99{,}5 & 99{,}9\\
Pima & \textbf{76{,}7} & 76{,}1 & 76{,}6 & 70{,}8 & 76{,}0 & \textbf{76{,}7} & 76{,}3 & 76{,}5 & 76{,}2\\
Breast & 96{,}6 & 96{,}5 & 96{,}4 & 95{,}3 & 96{,}4 & \textbf{96{,}8} & 96{,}3 & 96{,}0 & 95{,}6\\
Magic & 79{,}0 & 85{,}5 & 85{,}6 & 78{,}3 & 85{,}0 & \textbf{85{,}7} & 84{,}8 & 84{,}4 & 84{,}3\\
Parkinsons & 85{,}4 & 89{,}9 & 91{,}4 & \textbf{93{,}7} & 89{,}4 & 90{,}6 & 86{,}2 & 89{,}3 & 85{,}5\\
Ionosphere & 86{,}2 & 93{,}4 & \textbf{93{,}8} & 86{,}4 & 91{,}1 & 88{,}0 & 89{,}7 & 87{,}7 & 87{,}0\\
Spambase & 93{,}0 & 93{,}6 & 93{,}4 & 88{,}8 & \textbf{94{,}6} & 93{,}9 & 93{,}4 & 92{,}6 & 82{,}1\\
Sonar & 75{,}2 & 84{,}3 & \textbf{84{,}5} & 81{,}2 & 82{,}0 & 81{,}9 & 80{,}2 & 79{,}8 & 76{,}2\\
\midrule
Hạng trung bình & 6{,}55 & 3{,}15 & 2{,}35 & 7{,}40 & 4{,}60 & 3{,}00 & 5{,}35 & 6{,}00 & 6{,}60\\
\bottomrule
\end{tabular}}
\end{table}


\paragraph{Kết quả tổng quát.} \Cref{tab:bench} cho độ chính xác trung bình và \cref{fig:bench-diff} biểu diễn chênh lệch so với SVM hạt nhân Gauss; độ lệch chuẩn và kích thước mô hình được cho ở Phụ lục~B. Ba phương pháp dẫn đầu về độ chính xác là LS-SVM hạt nhân Gauss, mạng nơ-ron và SVM hạt nhân Gauss, với hạng trung bình lần lượt $2{,}35$, $3{,}00$ và $3{,}15$. PWL-C-SVM cộng tính đứng giữa bảng (hạng $5{,}35$), trên SVM tuyến tính ($6{,}55$) và kNN ($7{,}40$).

Bức tranh rõ hơn khi tách các bộ dữ liệu thành hai nhóm. Ở nhóm thứ nhất gồm Haberman, Transfusion, Banknote, Pima và Breast, mọi phương pháp dựa trên SVM, kể cả SVM tuyến tính, chênh nhau không quá khoảng $1{,}5$ điểm phần trăm: một biên tuyến tính đã gần như đủ tốt, và tính phi tuyến không đem lại lợi ích đáng kể. Ở nhóm thứ hai gồm Magic, Parkinsons, Ionosphere, Sonar, nơi SVM tuyến tính kém SVM hạt nhân Gauss từ $4$ đến $9$ điểm, PWL-C-SVM cộng tính thu hẹp đáng kể khoảng cách trên ba trong bốn bộ: nó cải thiện so với SVM tuyến tính $5{,}8$ điểm trên Magic, $3{,}5$ điểm trên Ionosphere và $5{,}0$ điểm trên Sonar, nhưng vẫn kém SVM hạt nhân Gauss từ $0{,}7$ đến $4{,}1$ điểm. Magic là trường hợp đáng chú ý nhất: PWL-C-SVM đạt $84{,}8\%$ so với $85{,}5\%$ của SVM hạt nhân Gauss, trong khi mô hình chỉ gồm $211$ số thực so với $7.694$ số, nhỏ hơn khoảng $36$ lần.

\begin{figure}[t]
\centering
\includegraphics[width=\linewidth]{fig_benchmark_diff.pdf}
\caption[Chênh lệch độ chính xác so với SVM hạt nhân Gauss]{Chênh lệch độ chính xác trung bình so với SVM hạt nhân Gauss (điểm phần trăm; giá trị âm nghĩa là kém hơn) của bốn phương pháp trên mười bộ dữ liệu.}
\label{fig:bench-diff}
\end{figure}

\begin{table}[t]
\centering
\caption[So sánh theo cặp với PWL-C-SVM cộng tính]{So sánh theo cặp giữa PWL-C-SVM cộng tính và từng phương pháp khác trên mười bộ dữ liệu: số bộ mà PWL-C-SVM cao hơn / bằng / thấp hơn (theo độ chính xác trung bình làm tròn đến $0{,}1$ điểm phần trăm), chênh lệch trung bình, và giá trị $p$ của kiểm định dấu hạng Wilcoxon hai phía.}
\label{tab:wilcoxon}
\small
\renewcommand{\arraystretch}{1.18}
\begin{tabular}{@{}l c r r@{}}
\toprule
\textbf{Phương pháp so sánh} & \textbf{Cao hơn / bằng / thấp hơn} & \textbf{Chênh lệch TB} & $p$\\
\midrule
SVM tuyến tính & 7 / 0 / 3 & $+1{,}66$ & 0{,}037\\
SVM-RBF & 1 / 0 / 9 & $-1{,}39$ & 0{,}014\\
LS-SVM-RBF & 0 / 1 / 9 & $-1{,}65$ & 0{,}008\\
kNN ($k=1$) & 8 / 0 / 2 & $+2{,}79$ & 0{,}105\\
Ik-SVM & 2 / 1 / 7 & $-0{,}68$ & 0{,}110\\
MLP-ReLU & 2 / 1 / 7 & $-0{,}82$ & 0{,}154\\
PWL-LS-SVM cộng tính & 7 / 0 / 3 & $+0{,}05$ & 0{,}432\\
PWL-C-SVM HH & 7 / 2 / 1 & $+1{,}97$ & 0{,}025\\
\bottomrule
\end{tabular}
\end{table}


\paragraph{Kiểm định thống kê.} \Cref{tab:wilcoxon} so sánh PWL-C-SVM cộng tính với từng phương pháp bằng kiểm định dấu hạng Wilcoxon trên mười cặp giá trị trung bình, theo khuyến nghị của Demšar~\cite{demsar2006}. Ở mức ý nghĩa $5\%$, PWL-C-SVM cộng tính kém SVM hạt nhân Gauss và LS-SVM hạt nhân Gauss một cách có ý nghĩa ($p=0{,}014$ và $p=0{,}008$). Nó cao hơn SVM tuyến tính trên bảy trong mười bộ, với chênh lệch trung bình $1{,}7$ điểm, và khác biệt này có ý nghĩa thống kê ($p=0{,}037$); khác biệt so với mạng nơ-ron và Ik-SVM không có ý nghĩa thống kê. Với chỉ mười bộ dữ liệu, các kiểm định này có lực thấp, nên “không bác bỏ” không có nghĩa là “tương đương”.

\paragraph{Ánh xạ HH và quy tắc chọn tham số.} PWL-C-SVM với ánh xạ HH có hạng trung bình thấp ($6{,}60$), và trên Spambase kết quả của nó rất bất ổn: ở lần chia đầu tiên, mô hình thất bại hoàn toàn ($40{,}0\%$, thấp hơn cả tỉ lệ $60{,}6\%$ của lớp đa số), trong khi ở bốn lần chia còn lại nó đạt từ $91{,}0\%$ đến $93{,}7\%$. Nguyên nhân nằm ở quy tắc sinh tham số chứ không ở mô hình. Dữ liệu Spambase rất thưa: sau khi chuẩn hóa, gần $78\%$ các giá trị bằng $0$, và phần lớn dữ liệu nằm sát một đỉnh của hình hộp $[0,1]^{57}$. Các siêu phẳng đi qua những điểm lấy đều trong hình hộp vì vậy thường không cắt qua vùng dữ liệu: ở mỗi lần chia, từ 18\% đến 22\% số đặc trưng bản lề bằng $0$ trên mọi điểm huấn luyện, và từ 17\% đến 22\% luôn hoạt động, tức là chỉ tuyến tính trên dữ liệu. Nghiêm trọng hơn, việc cố định $|p(1)|=1$ (\cref{rem:scale}) làm độ lớn của các đặc trưng phụ thuộc mạnh vào lần sinh ngẫu nhiên: giá trị lớn nhất của các đặc trưng trên tập huấn luyện dao động từ $423$ đến $3{,}3\cdot10^{5}$ giữa các lần chia. Ở lần chia đầu tiên, một siêu phẳng gần song song với trục $x(1)$ có hệ số lên tới $5{,}3\cdot10^{4}$, khiến bài toán tối ưu bị điều kiện hóa rất kém và bộ giải LIBLINEAR không hội tụ tới một nghiệm có ý nghĩa. Chỉ cần chuẩn hóa $\|p\|_2=1$, các đặc trưng trên dữ liệu huấn luyện không vượt quá $2{,}0$ ở mọi lần chia, và với cùng giao thức, PWL-C-SVM HH đạt $93{,}4\pm0{,}2\%$ trên Spambase. Nhìn chung, mạng nơ-ron với cùng số nơ-ron ẩn tốt hơn PWL-C-SVM HH trên tám trong mười bộ: học lớp ẩn từ dữ liệu có lợi hơn nhiều so với sinh nó ngẫu nhiên, và đó là cái giá của tính lồi đã được nêu ở Mục~\ref{sec:pwlsvm}.

\paragraph{So với bài báo gốc.} Kết quả trên không xác nhận nhận định của bài báo gốc rằng PWL-SVM có độ chính xác tương đương SVM hạt nhân Gauss~\cite{huang2013}. Nguyên nhân chính là bài báo chỉ dùng một lần chia dữ liệu: trên các bộ nhỏ như Parkinsons, với $98$ điểm kiểm tra, độ lệch chuẩn giữa các lần chia lên tới $2$--$4$ điểm phần trăm, nên một lần chia có thể cho kết quả rất lạc quan; bài báo báo cáo độ chính xác $100\%$ cho PWL-C-SVM trên bộ này, trong khi trung bình qua mười lần chia là $86{,}2\%$. Mặt khác, lợi thế về kích thước mô hình được xác nhận rõ ràng: với ánh xạ cộng tính, PWL-C-SVM cần từ $64$ đến $1.261$ số thực, nhỏ hơn SVM hạt nhân Gauss từ $2{,}5$ đến $36$ lần trên cả mười bộ dữ liệu (Phụ lục~B).

% ----------------------------------------------------------------------------
\subsection{Ảnh hưởng của số đặc trưng và của quy tắc chọn tham số}\label{sec:sensitivity}

Để trả lời câu hỏi CH3 một cách có kiểm soát, ta cố định mô hình là PWL-LS-SVM, giải chính xác bằng hồi quy ridge (\cref{prop:ridge}) để loại trừ ảnh hưởng của bộ giải, và thay đổi hai yếu tố: số đặc trưng trên mỗi chiều $M/n\in\{1,2,3,5,10,20\}$ và quy tắc chọn tham số. Bốn quy tắc được so sánh: ánh xạ cộng tính với lưới đều (quy tắc của bài báo gốc) và với lưới phân vị (\cref{rem:quantile}); ánh xạ HH với $|p(1)|=1$ (quy tắc của bài báo gốc) và với $\|p\|_2=1$ (\cref{rem:scale}). Thí nghiệm được thực hiện trên ba bộ dữ liệu có đặc điểm khác nhau --- Ionosphere ($n=33$), Magic ($n=10$, nhiều dữ liệu) và Spambase ($n=57$, rất thưa) --- với $5$ lần chia như ở Mục~\ref{sec:bench10} và $\gamma$ chọn bằng kiểm định chéo $5$ phần. Kết quả được vẽ ở \cref{fig:sensitivity}.

\begin{figure}[t]
\centering
\includegraphics[width=\linewidth]{fig_sensitivity.pdf}
\caption[Ảnh hưởng của số đặc trưng và quy tắc chọn tham số]{Độ chính xác trung bình (qua $5$ lần chia) của PWL-LS-SVM theo số đặc trưng trên mỗi chiều $M/n$, với bốn quy tắc chọn tham số. $M/n=1$ ứng với bộ phân loại tuyến tính.}
\label{fig:sensitivity}
\end{figure}

Có bốn nhận xét. Thứ nhất, độ chính xác tăng nhanh khi $M/n$ tăng từ $1$ (bộ phân loại tuyến tính) và bão hòa quanh $M/n=5$--$10$, phù hợp với quy tắc $M=10n$ của bài báo gốc; tăng tiếp lên $M/n=20$ hầu như không có lợi. Thứ hai, đặt nút theo phân vị cho kết quả tốt hơn hoặc ngang lưới đều trên cả ba bộ dữ liệu, và khác biệt rõ nhất trên những dữ liệu có phân bố lệch: trên Magic, lưới phân vị đạt $84{,}2\%$ ngay với $M/n=2$, trong khi lưới đều cần $M/n=10$ để đạt $84{,}3\%$; trên Spambase, lưới phân vị đạt $93{,}9\%$ so với $92{,}6\%$ của lưới đều ở $M/n=10$, đồng thời chỉ dùng $156$ đặc trưng thay vì $570$, vì các nút trùng nhau của những biến gần như luôn bằng $0$ bị loại bỏ. Thứ ba, ánh xạ HH cần nhiều đặc trưng hơn hẳn mới đạt tới độ chính xác của ánh xạ cộng tính: trên Magic, nó chỉ đuổi kịp lưới phân vị khi $M/n\ge10$ (và nhỉnh hơn $0{,}4$ điểm ở $M/n=20$), còn trên Spambase nó kém ở mọi giá trị của $M/n$; chỉ trên Ionosphere, nơi tương tác giữa các biến có vai trò, HH với $\|p\|_2=1$ đạt kết quả tốt nhất ($90{,}0\%$ ở $M/n=10$). Thứ tư, chuẩn hóa $\|p\|_2=1$ không bao giờ làm kết quả xấu đi đáng kể và cải thiện rõ trên Ionosphere ($90{,}0\%$ so với $86{,}9\%$ ở $M/n=10$); kết hợp với sự cố của LIBLINEAR trên Spambase, đây là một lý do mạnh để thay quy tắc $|p(1)|=1$ bằng $\|p\|_2=1$.

Từ đó rút ra một quy trình thực hành cho dữ liệu dạng bảng: bắt đầu với ánh xạ cộng tính, nút theo phân vị và $M/n$ khoảng $5$ đến $10$; chỉ dùng ánh xạ HH, với các siêu phẳng được chuẩn hóa, khi có dấu hiệu tương tác mạnh giữa các biến. Quy trình này được áp dụng cho bài toán tín dụng ở Mục~\ref{sec:application}.


% Generated by code/fill_numbers.py from chuong3_cost.tpl -- edit the template, not this file.
% ----------------------------------------------------------------------------
\section{Chi phí tính toán và bộ nhớ}\label{sec:cost}

Mục~\ref{sec:pwlsvm} dự báo rằng lợi thế của PWL-SVM nằm ở khâu dự đoán: mô hình chỉ gồm $w$, $w_0$ và tham số của ánh xạ, và việc dự đoán chỉ gồm các phép so sánh, nhân, cộng và lấy cực đại (\cref{tab:cost}). Mục này kiểm tra dự báo đó trên hai bộ dữ liệu lớn hơn: Magic với $9.510$ mẫu huấn luyện (chia đôi toàn bộ dữ liệu) và dữ liệu tín dụng sẽ được mô tả ở Mục~\ref{sec:application}, với $21.000$ mẫu huấn luyện. Để phép đo công bằng, mọi mô hình chạy trên một luồng xử lý của cùng một máy tính, khi không có tiến trình nặng nào khác; tham số được chọn trước bằng kiểm định chéo trên một mẫu con $4.000$ điểm, và chỉ thời gian của một lần huấn luyện với tham số đã chọn được đo. Thời gian dự đoán là trung vị của bảy lần dự đoán toàn bộ tập kiểm tra, chia cho số mẫu. PWL-SVM dùng ánh xạ cộng tính với nút phân vị và $M/n=10$, theo quy trình rút ra ở Mục~\ref{sec:sensitivity}; thẻ điểm là hồi quy logistic trên các biến rời rạc hóa, như ở Mục~\ref{sec:application}.


\paragraph{Dự đoán và bộ nhớ.} Kết quả ở \cref{tab:cost-measured} và \cref{fig:cost} xác nhận dự báo lý thuyết. PWL-LS-SVM cộng tính dự đoán mỗi mẫu trong $0{,}33$~$\mu$s trên Magic và $0{,}56$~$\mu$s trên dữ liệu tín dụng, nhanh hơn SVM hạt nhân Gauss khoảng 200 và 430 lần. Chênh lệch đến từ cấu trúc của mô hình: trên dữ liệu tín dụng, SVM hạt nhân Gauss lưu $239.111$ số thực, ứng với khoảng $8.854$ vectơ hỗ trợ, và mỗi lần dự đoán phải tính chừng ấy khoảng cách cùng hàm mũ; còn PWL-SVM cộng tính chỉ lưu $349$ số và thực hiện vài trăm phép so sánh, nhân và cộng, đúng như các công thức ở \cref{tab:cost}. Trên dữ liệu tín dụng, mạng nơ-ron dự đoán nhanh ngang PWL-SVM cộng tính, còn tổ hợp cây tăng cường chậm hơn khoảng 2{,}5 lần; cả hai mô hình cũng lớn hơn nhiều, với $6.605$ và khoảng $4.514$ số thực. Thẻ điểm có kích thước tương đương PWL-SVM cộng tính ($379$ so với $349$ số thực), và trong cài đặt dựa trên scikit-learn, nó dự đoán chậm hơn PWL-SVM khoảng 2{,}6 lần. Chỉ các mô hình tuyến tính trên biến gốc là nhỏ và nhanh hơn PWL-SVM, với cái giá là độ chính xác thấp hơn rõ rệt: trên Magic, hồi quy logistic đạt $78{,}7\%$ so với $84{,}8\%$ của PWL-LS-SVM cộng tính.

\begin{table}[t]
\centering
\caption[Chi phí huấn luyện, dự đoán và bộ nhớ]{Chi phí trên Magic ($9.510$ mẫu huấn luyện) và dữ liệu tín dụng ($21.000$ mẫu huấn luyện), đo trên một luồng xử lý: thời gian huấn luyện một mô hình với tham số đã chọn, thời gian dự đoán mỗi mẫu, số thực cần lưu, và độ chính xác (Magic) hoặc AUC (tín dụng) trên tập kiểm tra. Đây là một lần chia riêng, với tham số chọn trên mẫu con (và với Magic, tập huấn luyện lớn hơn nhiều), nên độ chính xác khác với ở \cref{tab:bench,tab:credit}.}
\label{tab:cost-measured}
\footnotesize
\setlength{\tabcolsep}{3pt}
\renewcommand{\arraystretch}{1.18}
\begin{tabular}{@{}l r r r c r r r c@{}}
\toprule
 & \multicolumn{4}{c}{\textbf{Magic}} & \multicolumn{4}{c}{\textbf{Tín dụng}}\\
\cmidrule(lr){2-5}\cmidrule(l){6-9}
\textbf{Mô hình} & \begin{tabular}[b]{@{}c@{}}\textbf{Huấn}\\\textbf{luyện (s)}\end{tabular} & \begin{tabular}[b]{@{}c@{}}\textbf{Dự đoán}\\\textbf{($\mu$s/mẫu)}\end{tabular} & \begin{tabular}[b]{@{}c@{}}\textbf{Số}\\\textbf{thực}\end{tabular} & \begin{tabular}[b]{@{}c@{}}\textbf{ĐCX}\\\textbf{(\%)}\end{tabular} & \begin{tabular}[b]{@{}c@{}}\textbf{Huấn}\\\textbf{luyện (s)}\end{tabular} & \begin{tabular}[b]{@{}c@{}}\textbf{Dự đoán}\\\textbf{($\mu$s/mẫu)}\end{tabular} & \begin{tabular}[b]{@{}c@{}}\textbf{Số}\\\textbf{thực}\end{tabular} & \textbf{AUC}\\
\midrule
SVM tuyến tính & 5{,}2 & 0{,}02 & 31 & 78{,}7 & 9{,}3 & 0{,}06 & 79 & 0{,}694\\
Hồi quy logistic & 0{,}00 & 0{,}02 & 31 & 78{,}7 & 0{,}02 & 0{,}05 & 79 & 0{,}716\\
Thẻ điểm & 0{,}02 & 0{,}69 & 191 & 84{,}6 & 0{,}07 & 1{,}4 & 379 & 0{,}768\\
PWL-C-SVM, cộng tính & 4{,}7 & 0{,}33 & 211 & 85{,}7 & 52 & 0{,}58 & 349 & 0{,}685\\
PWL-LS-SVM, cộng tính & 0{,}01 & 0{,}33 & 211 & 84{,}8 & 0{,}04 & 0{,}56 & 349 & 0{,}768\\
PWL-LS-SVM, HH & 0{,}01 & 0{,}40 & 1.111 & 85{,}1 & 0{,}07 & 1{,}3 & 6.631 & 0{,}740\\
SVM-RBF & 0{,}76 & 64{,}7 & 34.462 & 87{,}1 & 4{,}6 & 241 & 239.111 & 0{,}711\\
MLP-ReLU & 2{,}9 & 0{,}18 & 1.101 & 86{,}5 & 21 & 0{,}57 & 6.605 & 0{,}761\\
Tổ hợp cây tăng cường & 0{,}11 & 3{,}8 & 12.200 & 88{,}1 & 0{,}15 & 1{,}4 & 4.514 & 0{,}774\\
\bottomrule
\end{tabular}
\end{table}


\begin{figure}[t]
\centering
\includegraphics[width=\linewidth]{fig_cost.pdf}
\caption[Đánh đổi giữa độ chính xác và thời gian dự đoán]{Độ chính xác (Magic) hoặc AUC (dữ liệu tín dụng) theo thời gian dự đoán mỗi mẫu, trên thang logarit. Ô vuông là các biến thể PWL-SVM; điểm tròn là các phương pháp khác.}
\label{fig:cost}
\end{figure}

\paragraph{Huấn luyện.} PWL-LS-SVM là mô hình phi tuyến huấn luyện nhanh nhất: $0{,}01$ giây trên Magic và $0{,}04$ giây trên dữ liệu tín dụng, vì chỉ cần lập và giải một hệ tuyến tính cỡ $M\times M$ (\cref{prop:ridge}). PWL-C-SVM chậm hơn ($4{,}7$ và $52$ giây) vì bộ giải lặp của LIBLINEAR cần nhiều vòng để hội tụ. SVM hạt nhân Gauss cần $0{,}76$ và $4{,}6$ giây, và chi phí này tăng nhanh hơn tuyến tính theo số mẫu~\cite{chang2011}; mạng nơ-ron cần $2{,}9$ và $21$ giây. Trong một quy trình thực tế, nơi mô hình phải được huấn luyện lại nhiều lần để chọn tham số, khác biệt này còn được nhân lên theo kích thước lưới tham số.

Cần lưu ý rằng số đo tuyệt đối phụ thuộc vào cài đặt: PWL-SVM được cài bằng NumPy, còn LIBSVM là thư viện C đã được tối ưu nhiều năm. Trên một thiết bị nhúng, số phép tính và số thực cần lưu ở \cref{tab:cost} mới là thước đo phù hợp, và ở cả hai thước đo đó, PWL-SVM cộng tính có lợi thế lớn so với SVM hạt nhân Gauss.

% ----------------------------------------------------------------------------
\section{Ứng dụng: dự báo vỡ nợ thẻ tín dụng}\label{sec:application}

\subsection{Bài toán và dữ liệu}

Ngân hàng phát hành thẻ tín dụng cần ước lượng khả năng một khách hàng không thanh toán được dư nợ trong kỳ tới, để điều chỉnh hạn mức, ưu tiên nhắc nợ hoặc trích lập dự phòng rủi ro. Đây là một bài toán phân loại nhị phân điển hình của chấm điểm tín dụng~\cite{thomas2002}, và cả ba yêu cầu nêu ở Mở đầu đều có mặt: độ chính xác quyết định trực tiếp chi phí rủi ro; số hồ sơ lớn và việc chấm điểm được lặp lại định kỳ, nên chi phí dự đoán đáng kể; và tổ chức tín dụng cần giải thích được quyết định của mình với khách hàng cũng như với cơ quan quản lý.

Đề án dùng dữ liệu của Yeh và Lien~\cite{yeh2009} về $30.000$ chủ thẻ tín dụng của một ngân hàng ở Đài Loan, trong giai đoạn từ tháng 4 đến tháng 9 năm 2005. Nhãn cho biết khách hàng có vỡ nợ trong tháng tiếp theo hay không; có $6.636$ khách hàng vỡ nợ, tức $22{,}1\%$. Hai mươi ba biến giải thích thuộc bốn nhóm: hạn mức tín dụng, thông tin nhân khẩu học, lịch sử trả nợ sáu tháng gần nhất và các số tiền sao kê, thanh toán tương ứng (\cref{tab:credit-vars}).

\begin{table}[t]
\centering
\caption[Các biến của dữ liệu vỡ nợ thẻ tín dụng]{Các biến của dữ liệu vỡ nợ thẻ tín dụng~\cite{yeh2009} và cách xử lý trong đề án. Số tiền tính bằng Đài tệ.}
\label{tab:credit-vars}
\small
\renewcommand{\arraystretch}{1.2}
\begin{tabularx}{\linewidth}{@{}l X X@{}}
\toprule
\textbf{Biến} & \textbf{Ý nghĩa} & \textbf{Xử lý}\\
\midrule
LIMIT\_BAL & hạn mức tín dụng của khách hàng và gia đình & giữ nguyên\\
SEX & giới tính (1 nam, 2 nữ) & một biến chỉ báo “nữ”\\
EDUCATION & học vấn (1 sau đại học, 2 đại học, 3 trung học, 4 khác) & ba biến chỉ báo; các mã $0,5,6$ không được mô tả, gộp vào “khác”\\
MARRIAGE & hôn nhân (1 đã kết hôn, 2 độc thân, 3 khác) & hai biến chỉ báo; mã $0$ gộp vào “khác”\\
AGE & tuổi & giữ nguyên\\
PAY\_0, PAY\_2, \dots, PAY\_6 & trạng thái trả nợ các tháng 9, 8, \dots, 4: $-1$ trả đúng hạn, $k\ge1$ chậm $k$ tháng; các mã $-2$ và $0$ cũng xuất hiện & giữ nguyên như biến thứ bậc\\
BILL\_AMT1, \dots, BILL\_AMT6 & dư nợ sao kê các tháng 9, \dots, 4 & giữ nguyên\\
PAY\_AMT1, \dots, PAY\_AMT6 & số tiền đã thanh toán các tháng 9, \dots, 4 & giữ nguyên\\
\bottomrule
\end{tabularx}
\end{table}

Sau khi mã hóa các biến định danh, mỗi khách hàng được mô tả bởi $n=26$ đặc trưng số. Hai đặc điểm của dữ liệu quyết định thiết kế mô hình (\cref{fig:credit-overview}). Thứ nhất, quan hệ giữa lịch sử trả nợ và rủi ro là phi tuyến rõ rệt: tỉ lệ vỡ nợ gần như không đổi, khoảng $13$--$17\%$, với các trạng thái $-2$, $-1$ và $0$, rồi tăng vọt lên $34\%$ khi chậm một tháng và gần $70\%$ khi chậm hai tháng. Một mô hình tuyến tính theo PAY\_0 không thể mô tả bước nhảy này, trong khi một hàm tuyến tính từng phần với vài điểm gãy thì có thể. Thứ hai, các biến tiền tệ lệch rất mạnh: hạn mức trải từ $10$ nghìn đến $1$ triệu Đài tệ, nhưng một nửa số khách hàng có hạn mức không vượt quá $140$ nghìn. Lưới nút đều của bài báo gốc vì vậy đặt phần lớn các nút ở vùng gần như không có dữ liệu, còn lưới theo phân vị (\cref{rem:quantile}) đặt nút theo mật độ dữ liệu.

\begin{figure}[t]
\centering
\includegraphics[width=\linewidth]{fig_credit_overview.pdf}
\caption[Hai đặc điểm của dữ liệu tín dụng]{(a) Tỉ lệ vỡ nợ theo trạng thái trả nợ tháng 9 (PAY\_0); cột nhạt ứng với các nhóm có ít hơn $100$ khách hàng. (b) Phân bố hạn mức tín dụng cùng vị trí $9$ nút của lưới đều (nét đứt) và của lưới theo phân vị (nét liền).}
\label{fig:credit-overview}
\end{figure}

\subsection{Thiết kế thực nghiệm}

Dữ liệu được chia ngẫu nhiên có phân tầng theo tỉ lệ $70/30$, thành $21.000$ khách hàng huấn luyện và $9.000$ khách hàng kiểm tra, lặp lại $5$ lần. Độ đo chính là diện tích dưới đường cong ROC (AUC)~\cite{fawcett2006}: xác suất để một khách hàng vỡ nợ chọn ngẫu nhiên được mô hình chấm điểm rủi ro cao hơn một khách hàng không vỡ nợ chọn ngẫu nhiên. AUC không phụ thuộc vào ngưỡng quyết định và là độ đo chuẩn trong chấm điểm tín dụng~\cite{lessmann2015}. Để đánh giá thêm chất lượng phân loại, đề án chọn một ngưỡng trên tập huấn luyện sao cho tỉ lệ khách hàng bị dự báo vỡ nợ bằng tỉ lệ vỡ nợ quan sát được, rồi tính độ chính xác cân bằng (trung bình của độ nhạy và độ đặc hiệu) và độ đo F1 của lớp vỡ nợ trên tập kiểm tra. Quy tắc chọn ngưỡng này tương ứng với chính sách “đưa vào diện theo dõi một tỉ lệ khách hàng cố định”.

Mười mô hình được so sánh. Ba mô hình tuyến tính đóng vai trò chuẩn so sánh của ngành: hồi quy logistic trên các biến gốc, SVM tuyến tính, và \emph{thẻ điểm} --- hồi quy logistic trên các biến đã được rời rạc hóa, theo cách xây dựng thẻ điểm truyền thống~\cite{thomas2002}, trong đó mỗi biến liên tục được chia thành mười khoảng theo thập phân vị và mỗi giá trị của các biến PAY là một hạng mục riêng. Bốn biến thể PWL-SVM được thử: PWL-LS-SVM cộng tính với lưới đều (quy tắc của bài báo gốc) và với lưới phân vị, PWL-C-SVM cộng tính với lưới phân vị, và PWL-LS-SVM với ánh xạ HH; mọi biến thể dùng $M/n=10$. Ba mô hình phi tuyến làm đối chứng: SVM hạt nhân Gauss, mạng nơ-ron một lớp ẩn với $9n=234$ nơ-ron ReLU, huấn luyện bằng L-BFGS~\cite{liu1989}, và tổ hợp cây tăng cường~\cite{friedman2001} với thiết lập mặc định của scikit-learn, là một đại diện cho nhóm mô hình thường chính xác nhất trên dữ liệu dạng bảng~\cite{lessmann2015}. Tham số của mỗi mô hình được chọn bằng kiểm định chéo $5$ phần trên tập huấn luyện với tiêu chí AUC; riêng SVM hạt nhân Gauss và mạng nơ-ron được tinh chỉnh trên một mẫu con phân tầng gồm $5.000$ và $7.000$ khách hàng với kiểm định chéo $3$ phần, rồi huấn luyện lại trên toàn bộ tập huấn luyện, vì chi phí huấn luyện của chúng lớn hơn nhiều. Các lưới tham số là: $\gamma\in\{2^{-9},2^{-7},\dots,2^{7}\}$ cho SVM tuyến tính và mọi biến thể PWL-SVM; $\gamma\in\{2^{-2},2^{0},2^{2},2^{4}\}$ và $1/\sigma^2\in\{2^{-7},2^{-5},2^{-3},2^{-1}\}$ cho SVM hạt nhân Gauss; hệ số phạt $L_2$ trong $\{10^{-2},10^{-1},1,10\}$ cho mạng nơ-ron; nghịch đảo của hệ số phạt trong $\{10^{-3},10^{-2},\dots,10\}$ cho hồi quy logistic và thẻ điểm. Tổ hợp cây tăng cường dùng thiết lập mặc định, có cơ chế dừng sớm trên một phần dữ liệu kiểm định nội bộ.


% Generated by code/fill_numbers.py from chuong3_credit_b.tpl -- edit the template, not this file.
\subsection{Kết quả}\label{sec:credit-results}

\begin{table}[t]
\centering
\caption[Kết quả trên dữ liệu vỡ nợ thẻ tín dụng]{Kết quả trên dữ liệu vỡ nợ thẻ tín dụng (trung bình $\pm$ độ lệch chuẩn qua $5$ lần chia $70/30$). AUC: diện tích dưới đường ROC; ĐCX-CB: độ chính xác cân bằng; F1 tính cho lớp vỡ nợ; ngưỡng được chọn trên tập huấn luyện sao cho tỉ lệ dự báo vỡ nợ bằng tỉ lệ vỡ nợ quan sát. Cột cuối: số thực cần lưu (trung vị); thời gian dự đoán được đo riêng ở \cref{sec:cost}.}
\label{tab:credit}
\footnotesize
\setlength{\tabcolsep}{3pt}
\renewcommand{\arraystretch}{1.18}
\begin{tabular}{@{}l c c c r@{}}
\toprule
\textbf{Mô hình} & \textbf{AUC} & \textbf{ĐCX-CB (\%)} & \textbf{F1 (\%)} & \textbf{Số thực}\\
\midrule
Hồi quy logistic & $0{,}724\pm0{,}009$ & $68{,}7\pm0{,}6$ & $51{,}2\pm0{,}9$ & 79\\
Thẻ điểm (logistic, biến rời rạc hóa) & $0{,}773\pm0{,}006$ & $70{,}3\pm0{,}6$ & $53{,}7\pm1{,}0$ & 381\\
SVM tuyến tính & $0{,}706\pm0{,}008$ & $68{,}5\pm0{,}8$ & $50{,}9\pm1{,}4$ & 79\\
\midrule
PWL-LS-SVM, cộng tính, lưới đều & $0{,}772\pm0{,}006$ & $70{,}3\pm0{,}5$ & $53{,}6\pm0{,}9$ & 547\\
PWL-LS-SVM, cộng tính, phân vị & $0{,}776\pm0{,}006$ & $70{,}4\pm0{,}7$ & $53{,}8\pm1{,}1$ & 349\\
\quad và ràng buộc đơn điệu (Mục~\ref{sec:monotone}) & $0{,}775\pm0{,}006$ & $70{,}3\pm0{,}8$ & $53{,}7\pm1{,}2$ & 349\\
PWL-C-SVM, cộng tính, phân vị & $0{,}744\pm0{,}007$ & $67{,}8\pm0{,}8$ & $49{,}9\pm1{,}4$ & 349\\
PWL-LS-SVM, HH & $0{,}743\pm0{,}007$ & $69{,}5\pm0{,}6$ & $52{,}4\pm1{,}0$ & 6.631\\
\midrule
SVM-RBF & $0{,}716\pm0{,}010$ & $69{,}2\pm0{,}6$ & $52{,}0\pm1{,}0$ & 243.134\\
MLP-ReLU & $0{,}768\pm0{,}005$ & $70{,}0\pm0{,}4$ & $53{,}3\pm0{,}7$ & 6.605\\
Tổ hợp cây tăng cường & $0{,}779\pm0{,}006$ & $70{,}3\pm0{,}7$ & $53{,}7\pm1{,}1$ & 4.636\\
\bottomrule
\end{tabular}
\end{table}


\begin{figure}[t]
\centering
\includegraphics[width=0.62\linewidth]{fig_credit_roc.pdf}
\caption[Đường cong ROC trên dữ liệu tín dụng]{Đường cong ROC trên tập kiểm tra của lần chia đầu tiên cho bốn mô hình; con số trong ngoặc là AUC trung bình qua năm lần chia.}
\label{fig:credit-roc}
\end{figure}

\Cref{tab:credit} và \cref{fig:credit-roc} cho thấy các mô hình chia thành hai nhóm rõ rệt. Nhóm dẫn đầu gồm tổ hợp cây tăng cường (AUC trung bình 0{,}779), PWL-LS-SVM cộng tính với nút phân vị (0{,}776), thẻ điểm logistic trên các biến đã rời rạc hóa (0{,}773) và PWL-LS-SVM cộng tính với lưới đều (0{,}772). Nhóm còn lại gồm hồi quy logistic trên các biến gốc (0{,}724), SVM tuyến tính (0{,}706) và SVM hạt nhân Gauss (0{,}716); mạng nơ-ron (0{,}768) đứng giữa hai nhóm. Như vậy, AUC của PWL-LS-SVM cộng tính với nút phân vị chỉ kém tổ hợp cây tăng cường một chút. Độ chính xác cân bằng và độ đo F1 ở ngưỡng đã chọn dẫn tới cùng kết luận: bốn mô hình của nhóm dẫn đầu gần như ngang nhau và cao hơn rõ rệt hồi quy logistic, SVM tuyến tính và SVM hạt nhân Gauss. Kết quả rất ổn định giữa các lần chia: độ lệch chuẩn của AUC chỉ khoảng 0{,}006.

Khoảng cách 0{,}052 đơn vị AUC giữa PWL-LS-SVM và hồi quy logistic trên các biến gốc cần được hiểu đúng. Phần lớn khoảng cách đó đến từ quan hệ phi tuyến giữa các trạng thái trả nợ và rủi ro (\cref{fig:credit-overview}a), mà một mô hình tuyến tính theo PAY\_0 không mô tả được. Trong thực tế, các thẻ điểm không dùng trực tiếp biến gốc mà rời rạc hóa từng biến thành các khoảng rồi mới hồi quy logistic~\cite{thomas2002}, và thẻ điểm xây dựng theo cách đó ở \cref{tab:credit} đã thu hẹp phần lớn khoảng cách. Thẻ điểm cũng là đối chứng công bằng nhất cho PWL-SVM cộng tính: cả hai đều là mô hình cộng tính, đều dựa trên các khoảng của từng biến, chỉ khác ở chỗ đóng góp của mỗi biến là hàm bậc thang trong thẻ điểm và là đường gấp khúc liên tục trong PWL-SVM. So với thẻ điểm, PWL-LS-SVM cao hơn ở bốn trong năm lần chia, với chênh lệch trung bình 0{,}002 đơn vị AUC, trong khi hai mô hình có kích thước tương đương (349 và 381 số thực). Chênh lệch này nhỏ, nhưng trong chấm điểm tín dụng, ngay cả những cải thiện nhỏ về khả năng phân biệt cũng có thể mang lại lợi ích kinh tế đáng kể khi áp dụng cho một danh mục lớn~\cite{lessmann2015}.

Ba quan sát khác đáng chú ý. Thứ nhất, lưới phân vị tốt hơn lưới đều (0{,}776 so với 0{,}772) với mô hình nhỏ hơn (349 so với 547 số thực), đúng như phân tích ở \cref{fig:credit-overview}b. Thứ hai, với cùng ánh xạ, PWL-C-SVM cho AUC thấp hơn hẳn PWL-LS-SVM (0{,}744). Điều này có thể giải thích bằng hàm mất mát: mất mát bản lề chỉ quan tâm đến các điểm nằm gần hoặc vi phạm lề, nên giá trị hàm quyết định ở xa biên không mang nhiều thông tin về mức độ rủi ro; còn mất mát bình phương của LS-SVM, tương đương hồi quy trên nhãn $\pm1$ (\cref{prop:ridge}), buộc điểm số của mọi khách hàng phải phản ánh nhãn của họ, và vì thế cho một thứ hạng tốt hơn. SVM hạt nhân Gauss, cũng dùng mất mát bản lề, gặp cùng hạn chế. Thứ ba, ánh xạ HH (0{,}743) kém ánh xạ cộng tính: trên dữ liệu này, ảnh hưởng của từng biến quan trọng hơn tương tác giữa các biến, và các siêu phẳng ngẫu nhiên lãng phí khả năng biểu diễn, phù hợp với kết luận của Mục~\ref{sec:sensitivity}.

\subsection{Diễn giải mô hình}\label{sec:credit-interp}

Vì ánh xạ là cộng tính, hàm quyết định của PWL-LS-SVM tách thành tổng các đóng góp của từng biến,
\begin{equation}\label{eq:additive-decomp}
f(x)=w_0+\sum_{i=1}^{n}h_i\bigl(x(i)\bigr),
\end{equation}
trong đó mỗi $h_i$ là một hàm tuyến tính từng phần một biến, có không quá chín điểm gãy đặt tại các phân vị của biến đó (\cref{prop:1d}). Mỗi $h_i$ có thể được vẽ ra và đọc trực tiếp, giống như một thẻ điểm trong đó số điểm cộng thêm cho mỗi khách hàng là một đường gấp khúc theo giá trị của biến. \Cref{fig:credit-shapes} vẽ sáu hàm $h_i$ có ảnh hưởng lớn nhất của mô hình huấn luyện trên lần chia đầu tiên, sau khi trừ đi giá trị trung bình trên tập huấn luyện; giá trị dương nghĩa là rủi ro cao hơn một khách hàng trung bình.

\begin{figure}[t]
\centering
\includegraphics[width=\linewidth]{fig_credit_shapes.pdf}
\caption[Các hàm đóng góp của mô hình PWL-SVM cộng tính]{Sáu hàm đóng góp $h_i$ có ảnh hưởng lớn nhất của PWL-LS-SVM cộng tính (nút phân vị) trên dữ liệu tín dụng, đã trừ giá trị trung bình. Nét đậm: mô hình của lần chia đầu tiên, nét chấm dọc là vị trí các nút của nó; nét mảnh: mô hình của bốn lần chia còn lại. Giá trị dương ứng với rủi ro vỡ nợ cao hơn.}
\label{fig:credit-shapes}
\end{figure}

Hàm đóng góp của PAY\_0 thể hiện rõ cấu trúc mà hồi quy logistic trên biến gốc bỏ lỡ (so sánh với \cref{fig:credit-overview}a). Ba mã không chậm trả ($-2$, $-1$ và $0$) có đóng góp thấp và gần nhau; mã $-1$ cao hơn hai mã còn lại một chút, đúng như tỉ lệ vỡ nợ quan sát được (16{,}8\% so với 13{,}2\% và 12{,}8\%). Đóng góp tăng mạnh khi khách hàng chậm trả một tháng (tỉ lệ vỡ nợ 33{,}9\%) và đạt cực đại khi chậm hai tháng (69{,}1\%); chênh lệch giữa mã $0$ và mã $2$ là $0{,}86$ đơn vị điểm, lớn hơn toàn bộ biên độ của bất kỳ biến nào khác. Đóng góp của hạn mức tín dụng giảm dần theo hạn mức --- điều hợp lý, vì hạn mức do chính ngân hàng xác định dựa trên đánh giá tín nhiệm --- và khoảng 60\% mức giảm diễn ra ở vùng hạn mức dưới $100$ nghìn Đài tệ. Trạng thái trả nợ của các tháng trước (PAY\_2, PAY\_3, PAY\_4) có hình dạng tương tự PAY\_0 nhưng biên độ nhỏ hơn nhiều: thông tin gần nhất là quan trọng nhất.

Hai đặc điểm khác của các hàm đóng góp cần được đọc thận trọng. Thứ nhất, sau mức chậm hai tháng, đóng góp của PAY\_0 và nhất là của PAY\_4 lại giảm, trái với trực giác. Đây không phải là quy luật của dữ liệu mà là hệ quả của việc có quá ít dữ liệu: chỉ 1{,}5\% khách hàng có PAY\_0 từ $3$ trở lên, tỉ lệ vỡ nợ của nhóm này (71{,}9\%) không thấp hơn của nhóm chậm hai tháng, và những khách hàng này thường chậm trả nhiều tháng liền, nên mô hình có thể phân bổ rủi ro giữa các biến PAY theo nhiều cách gần tương đương. Hệ số góc của đoạn cuối, được ước lượng từ rất ít điểm, vì vậy không đáng tin cậy; Mục~\ref{sec:monotone} sẽ khắc phục hiện tượng này. Thứ hai, các biến dư nợ sao kê BILL\_AMT1--6 tương quan rất mạnh với nhau (hệ số tương quan giữa hai tháng liền kề từ $0{,}92$ đến $0{,}95$), nên mô hình có thể chuyển ảnh hưởng qua lại giữa chúng: chẳng hạn, đóng góp của BILL\_AMT3 tăng theo dư nợ, còn đóng góp của BILL\_AMT4 lại giảm. Với những nhóm biến như vậy, chỉ tổng đóng góp của cả nhóm mới nên được diễn giải. Đây là hạn chế chung của mọi mô hình cộng tính khi các biến tương quan mạnh, không riêng gì PWL-SVM.

Các nét mảnh ở \cref{fig:credit-shapes} cho thấy các hàm đóng góp thay đổi thế nào khi dữ liệu huấn luyện thay đổi. Ở vùng có nhiều dữ liệu, hình dạng gần như không đổi giữa các lần chia: giữa phân vị $5\%$ và $95\%$ của hạn mức tín dụng, đóng góp của LIMIT\_BAL ở bốn lần chia còn lại lệch khỏi lần chia đầu tiên không quá 15\% biên độ của nó, và ở các trạng thái từ $-2$ đến $2$, đóng góp của PAY\_0 lệch không quá 10\% biên độ. Thứ hạng của các biến cũng ổn định: ở cả năm lần chia, PAY\_0 là biến có ảnh hưởng lớn nhất và LIMIT\_BAL đứng thứ hai; thứ tự của các biến tiếp theo, vốn có mức ảnh hưởng gần nhau, thì thay đổi giữa các lần chia. Ngược lại, ở các trạng thái chậm trả từ ba tháng trở lên, các lần chia cho những đường rất khác nhau: độ lệch chuẩn giữa các lần chia của đóng góp, tính trung bình trên sáu biến PAY, ở những trạng thái này lớn gấp khoảng 8 lần so với ở các trạng thái từ $-2$ đến $2$ --- thêm một bằng chứng rằng các đoạn đó không đáng tin cậy.

\begin{figure}[t]
\centering
\includegraphics[width=\linewidth]{fig_credit_explain.pdf}
\caption[Mức độ ảnh hưởng của các biến và giải thích một hồ sơ]{(a) Mức độ ảnh hưởng của mười biến quan trọng nhất, đo bằng độ lệch chuẩn của đóng góp $h_i$ trên tập huấn luyện; với các biến định danh, đó là độ lệch chuẩn của tổng đóng góp của các biến chỉ báo tương ứng. (b) Tám đóng góp lớn nhất, so với khách hàng trung bình, của hồ sơ có điểm rủi ro cao nhất trong số các khách hàng vỡ nợ ở tập kiểm tra.}
\label{fig:credit-explain}
\end{figure}

\Cref{fig:credit-explain}a xếp hạng các biến theo mức độ ảnh hưởng. PAY\_0 vượt xa mọi biến khác; tiếp theo, với mức ảnh hưởng nhỏ hơn nhiều, là hạn mức tín dụng, các trạng thái trả nợ của những tháng trước và các số tiền sao kê; các biến nhân khẩu học (học vấn, hôn nhân, giới tính, tuổi) có ảnh hưởng không đáng kể. Vì hàm quyết định~\eqref{eq:additive-decomp} là một tổng, mỗi quyết định của mô hình được giải thích một cách chính xác chứ không phải xấp xỉ: điểm của một khách hàng bằng điểm trung bình trên tập huấn luyện cộng với tổng các đóng góp của từng biến so với khách hàng trung bình. \Cref{fig:credit-explain}b minh họa điều đó với khách hàng có điểm rủi ro cao nhất trong số những người vỡ nợ ở tập kiểm tra. Điểm của khách hàng này là $0{,}96$, trong khi điểm trung bình là $-0{,}56$ và ngưỡng phân loại là $-0{,}44$. Trong phần chênh lệch $1{,}51$ so với điểm trung bình, riêng việc chậm trả hai tháng ở tháng 9 đóng góp $0{,}71$, lịch sử chậm trả ở năm tháng trước đó đóng góp thêm $0{,}63$, và hạn mức thấp ($20.000$ Đài tệ) đóng góp $0{,}13$. Một lời giải thích như vậy có thể được trình bày trực tiếp cho cán bộ tín dụng hay cho khách hàng. Với SVM hạt nhân Gauss hay tổ hợp cây, việc giải thích một quyết định đòi hỏi những kỹ thuật bổ sung, chẳng hạn giá trị SHAP~\cite{lundberg2017}, được tính sau khi huấn luyện và phức tạp hơn nhiều so với việc đọc trực tiếp các hàm $h_i$.

\subsection{Ràng buộc đơn điệu}\label{sec:monotone}

Người làm tín dụng kỳ vọng rằng rủi ro không giảm khi số tháng chậm trả tăng, và không tăng khi hạn mức tín dụng tăng. Mô hình không ràng buộc chỉ thỏa mãn một phần kỳ vọng này: ở cả năm lần chia, ít nhất một trong sáu hàm $h_i$ của các biến PAY có một đoạn với hệ số góc âm ở vùng từ trạng thái $0$ trở đi, và hàm của LIMIT\_BAL có đoạn đi lên ở cả năm lần chia. Như đã phân tích ở trên, đó là những đoạn được ước lượng từ rất ít dữ liệu, không phải là quy luật thật. Một mô hình trái với tri thức nghiệp vụ khó được chấp nhận khi thẩm định, dù độ chính xác của nó cao; vì vậy ta đưa kỳ vọng đơn điệu vào chính bài toán huấn luyện.

Với PWL-SVM cộng tính, điều này có thể làm mà vẫn giữ tính lồi. Gọi $t_{i1}<\dots<t_{iK}$ là các nút của biến thứ $i$ và $s_{i0},\dots,s_{iK}$ là hệ số góc của $h_i$ trên $K+1$ đoạn mà các nút chia ra. Khi đó
\begin{equation}\label{eq:slope-param}
h_i(t)=\sum_{j=0}^{K}s_{ij}\,G_{ij}(t),\qquad
G_{ij}(t)=\min\bigl\{(t-t_{ij})_+,\ t_{i,j+1}-t_{ij}\bigr\},
\end{equation}
với quy ước $t_{i0}=-\infty$, $t_{i,K+1}=+\infty$ (tức $G_{i0}(t)=\min\{t,t_{i1}\}$ và $G_{iK}(t)=(t-t_{iK})_+$): $G_{ij}(t)$ là độ dài phần đoạn thứ $j$ nằm bên trái $t$. Theo \cref{prop:1d}, hệ số của $x(i)$ trong ánh xạ~\eqref{eq:feature-map} bằng $s_{i0}$ và hệ số của $\bigl(x(i)-t_{ij}\bigr)_+$ bằng $s_{ij}-s_{i,j-1}$. Vì vậy~\eqref{eq:slope-param} chỉ là một cách tham số hóa khác của cùng một mô hình, và bài toán~\eqref{eq:ridge} trở thành
\[
\min_{s,\,w_0}\ \frac12\sum_{i=1}^{n}\Bigl[s_{i0}^2+\sum_{j=1}^{K}\bigl(s_{ij}-s_{i,j-1}\bigr)^2\Bigr]+\frac\gamma2\sum_{k=1}^{N}\Bigl(y_k-w_0-\sum_{i=1}^{n}h_i\bigl(x_k(i)\bigr)\Bigr)^2 .
\]
Hàm $h_i$ không giảm khi và chỉ khi $s_{ij}\ge0$ với mọi $j$, và không tăng khi và chỉ khi $s_{ij}\le0$ với mọi $j$; muốn $h_i$ đơn điệu chỉ trên một phần của trục, chẳng hạn trên nửa đường thẳng $t\ge t_{ij_0}$, ta chỉ cần áp các bất đẳng thức đó cho những đoạn nằm trong phần ấy. Thêm các ràng buộc này, ta được một bài toán bình phương tối thiểu với ràng buộc cận trên các biến --- một bài toán lồi, giải được chính xác bằng thuật toán BVLS~\cite{stark1995} có sẵn trong SciPy~\cite{virtanen2020}. Cũng như ở \cref{rem:convex}, tri thức tiên nghiệm được đưa vào mô hình mà không phải từ bỏ tính lồi. Khi không có ràng buộc nào, lời giải trùng với lời giải của \cref{prop:ridge} (sai khác dưới $10^{-9}$), điều được dùng để kiểm tra cài đặt.

\begin{figure}[t]
\centering
\includegraphics[width=\linewidth]{fig_credit_monotone.pdf}
\caption[Hàm đóng góp khi có và không có ràng buộc đơn điệu]{Hàm đóng góp của PAY\_0, PAY\_4 và LIMIT\_BAL trong mô hình không ràng buộc (nét liền) và mô hình có ràng buộc đơn điệu (nét đứt), lần chia đầu tiên. Vùng tô xám: các trạng thái chậm trả từ ba tháng trở lên, chỉ chiếm chưa tới 2\% số khách hàng.}
\label{fig:credit-monotone}
\end{figure}

Ta ràng buộc sáu hàm của các biến PAY không giảm từ trạng thái $0$ trở đi --- ba mã $-2$, $-1$ và $0$ chỉ những tình trạng không chậm trả khác nhau, không có thứ tự tự nhiên, nên được để tự do --- và hàm của LIMIT\_BAL không tăng trên toàn miền; các biến khác, các nút và hằng số $\gamma$ đã chọn cho mô hình không ràng buộc được giữ nguyên. \Cref{fig:credit-monotone} cho thấy tác động của ràng buộc: sau mức chậm hai tháng, đóng góp của PAY\_0 và PAY\_4 giữ nguyên thay vì giảm, và đường của LIMIT\_BAL trở thành đơn điệu giảm, trong khi hình dạng ở vùng có nhiều dữ liệu hầu như không đổi. Độ chính xác hầu như không thay đổi: AUC trung bình là 0{,}775, so với 0{,}776 của mô hình không ràng buộc, và mô hình có ràng buộc vẫn cao hơn thẻ điểm (0{,}773) (\cref{tab:credit}). Như vậy, với ràng buộc đơn điệu, PWL-LS-SVM cộng tính cho một mô hình vừa chính xác vừa nhất quán với tri thức nghiệp vụ.

\subsection{Khả năng triển khai}

Mô hình PWL-LS-SVM cộng tính trên dữ liệu tín dụng gồm 349 số thực: các nút và hệ số của $20$ đường gấp khúc cho các biến số, hệ số của $6$ biến chỉ báo, hệ số chặn và hai vectơ chuẩn hóa. Để chấm điểm một hồ sơ, với mỗi biến ta chỉ cần xác định giá trị rơi vào đoạn nào giữa các nút rồi tính một hàm bậc nhất, tổng cộng vài trăm phép so sánh, nhân và cộng. Một mô hình như vậy có thể được cài đặt bằng vài dòng lệnh SQL trong kho dữ liệu của ngân hàng, bằng một bảng tính, hoặc trên thiết bị có tài nguyên rất hạn chế, mà không cần thư viện học máy nào; ràng buộc đơn điệu không làm thay đổi cấu trúc này. Việc huấn luyện lại khi có dữ liệu mới chỉ mất chưa tới một giây (\cref{tab:cost-measured}). Để so sánh, SVM hạt nhân Gauss cần lưu khoảng 243.134 số thực, và tổ hợp cây tăng cường gồm từ 34 đến 47 cây quyết định với khoảng 4.636 tham số; cả hai đều khó kiểm tra và giải thích hơn nhiều.

% Generated by code/fill_numbers.py from chuong3_end.tpl -- edit the template, not this file.
% ----------------------------------------------------------------------------
\section{Thảo luận}\label{sec:discussion}

\paragraph{Trả lời các câu hỏi nghiên cứu.} Các thực nghiệm của chương cho phép trả lời bốn câu hỏi nêu ở Mở đầu.

\emph{(CH1) Độ chính xác.} Trên dữ liệu hai chiều, PWL-SVM dạng HH đạt độ chính xác gần SVM hạt nhân Gauss, trừ trên bàn cờ (\cref{tab:exp2d}). Trên mười bộ dữ liệu chuẩn, PWL-C-SVM cộng tính đứng giữa SVM tuyến tính và SVM hạt nhân Gauss: nó kém SVM hạt nhân Gauss trung bình 1{,}4 điểm phần trăm, một khác biệt có ý nghĩa thống kê, và hơn SVM tuyến tính trung bình 1{,}7 điểm, với phần cải thiện tập trung ở những bộ dữ liệu phi tuyến rõ rệt (\cref{tab:bench,tab:wilcoxon}). Nhận định của bài báo gốc rằng PWL-SVM chính xác tương đương SVM hạt nhân Gauss vì vậy không được xác nhận khi thực nghiệm được lặp lại. Trên dữ liệu tín dụng, ngược lại, PWL-LS-SVM cộng tính thuộc nhóm mô hình tốt nhất, cùng với tổ hợp cây tăng cường và thẻ điểm logistic (\cref{tab:credit}).

\emph{(CH2) Chi phí.} Lợi thế về chi phí được xác nhận rõ ràng. Với ánh xạ cộng tính, mô hình nhỏ hơn SVM hạt nhân Gauss từ $2{,}5$ đến $36$ lần trên mười bộ dữ liệu chuẩn và khoảng 690 lần trên dữ liệu tín dụng; thời gian dự đoán nhanh hơn khoảng 200 lần trên Magic và 430 lần trên dữ liệu tín dụng (\cref{tab:cost-measured}). PWL-LS-SVM còn là mô hình phi tuyến huấn luyện nhanh nhất.

\emph{(CH3) Cấu hình.} Ánh xạ cộng tính là lựa chọn mặc định tốt cho dữ liệu dạng bảng: rẻ, dễ diễn giải và thường chính xác nhất trong các dạng ánh xạ. Ánh xạ HH chỉ có lợi khi tương tác giữa các biến đáng kể (bàn cờ, xoắn ốc, Ionosphere); kích thước của nó tăng gần như bậc hai theo số chiều, và quy tắc sinh tham số của bài báo gốc có thể thất bại trên dữ liệu thưa. Độ chính xác bão hòa quanh $M/n=5$--$10$. Hai điều chỉnh nhỏ của đề án --- nút theo phân vị và chuẩn hóa $\|p\|_2=1$ --- cải thiện kết quả và độ ổn định số mà không thay đổi bản chất của phương pháp.

\emph{(CH4) Ứng dụng.} Trên bài toán vỡ nợ thẻ tín dụng, PWL-LS-SVM cộng tính với nút phân vị đạt AUC 0{,}776, chỉ kém một chút so với tổ hợp cây tăng cường (0{,}779) và cao hơn thẻ điểm logistic truyền thống (0{,}773) ở bốn trong năm lần chia, trong khi mô hình chỉ gồm 349 số thực, dự đoán mỗi hồ sơ trong khoảng 0{,}56 micro giây và đọc được như một thẻ điểm tuyến tính từng phần (\cref{fig:credit-shapes}). Khi thêm ràng buộc đơn điệu để mô hình nhất quán với tri thức nghiệp vụ, AUC hầu như không đổi (0{,}775 so với 0{,}776) (Mục~\ref{sec:monotone}).

\paragraph{Khi nào nên dùng PWL-SVM.} Từ các kết quả trên, PWL-SVM cộng tính phù hợp với dữ liệu dạng bảng có số chiều vừa phải, khi mô hình cần nhỏ, dự đoán nhanh hoặc phải giải thích được: chấm điểm tín dụng, các hệ thống nhúng, hay bất cứ bối cảnh nào mà hồi quy logistic đang được dùng vì tính minh bạch. Khi độ chính xác là tiêu chí duy nhất và tài nguyên không bị hạn chế, SVM hạt nhân Gauss, mạng nơ-ron hay tổ hợp cây thường là lựa chọn tốt hơn. Với dữ liệu nhiều chiều, thưa hoặc có tương tác phức tạp, PWL-SVM với các siêu phẳng sinh ngẫu nhiên không phải là công cụ thích hợp; khi đó nên học cả các siêu phẳng, tức là chuyển sang mạng nơ-ron ReLU.

\paragraph{Giới hạn của thực nghiệm.} Kết luận của chương dựa trên mười bộ dữ liệu chuẩn và một bài toán ứng dụng, nên lực của các kiểm định thống kê còn hạn chế. Các lưới tham số được chọn theo khuyến nghị thông dụng và có thể chưa tối ưu cho mọi mô hình; SVM hạt nhân Gauss và mạng nơ-ron trên dữ liệu tín dụng được tinh chỉnh trên mẫu con; thẻ điểm đối chứng dùng một cách rời rạc hóa đơn giản (mười khoảng theo thập phân vị), trong khi các thẻ điểm thực tế thường được xây dựng với các khoảng chọn cẩn thận hơn. Bộ giải LIBLINEAR chính quy hóa nhẹ hệ số chặn, khác với bài toán~\eqref{eq:pwl-c-svm}. Dữ liệu tín dụng được thu thập ở Đài Loan năm 2005, nên các hàm đóng góp tìm được minh họa phương pháp chứ không nhất thiết phản ánh hành vi của khách hàng Việt Nam hiện nay. Cuối cùng, thời gian đo được phụ thuộc vào phần cứng và cách cài đặt.

% ----------------------------------------------------------------------------
\section*{Tiểu kết Chương 3}
\phantomsection\addcontentsline{toc}{section}{Tiểu kết Chương 3}

Chương 3 đã cài đặt PWL-SVM, kiểm chứng tính đúng đắn của cài đặt bằng các phép thử số, và đánh giá phương pháp qua ba lớp thực nghiệm: dữ liệu hai chiều, mười bộ dữ liệu chuẩn và một bài toán ứng dụng. Kết quả cho một bức tranh nhất quán: PWL-SVM không chính xác hơn các mô hình phi tuyến mạnh như SVM hạt nhân Gauss, nhưng nó cho những mô hình nhỏ hơn từ vài lần đến hàng trăm lần, dự đoán nhanh hơn nhiều, huấn luyện bằng một bài toán lồi và, với ánh xạ cộng tính, có thể đọc được từng thành phần. Trên bài toán dự báo vỡ nợ thẻ tín dụng, nơi các tính chất đó được coi trọng, PWL-LS-SVM cộng tính với nút phân vị đạt độ chính xác ngang những mô hình tốt nhất, và với ràng buộc đơn điệu, cho một mô hình nhất quán với tri thức nghiệp vụ mà không mất độ chính xác.

